% !TeX program = xelatex
% =========================================================
% ElegantBook 数学知识讲义模板
% =========================================================

\documentclass[
    lang=cn,
    color=blue,
    mode=fancy,
    thmcnt=section,
    scheme=chinese
]{elegantbook}


% =========================================================
% 数学宏包
% =========================================================

\usepackage{amsmath}
\usepackage{amssymb}
\usepackage{mathtools}
\usepackage{bm}

% 图形
\usepackage{tikz}
\usepackage{pgfplots}
\pgfplotsset{compat=1.18}

% 表格
\usepackage{booktabs}
\usepackage{array}

% 多栏
\usepackage{multicol}


% =========================================================
% 常用数学命令
% =========================================================

% 数集
\newcommand{\R}{\mathbb{R}}
\newcommand{\C}{\mathbb{C}}
\newcommand{\N}{\mathbb{N}}
\newcommand{\Z}{\mathbb{Z}}
\newcommand{\Q}{\mathbb{Q}}

% 微分与常数
\newcommand{\dd}{\,\mathrm{d}}
\newcommand{\e}{\mathrm{e}}
\newcommand{\ii}{\mathrm{i}}

% 绝对值与范数
\newcommand{\abs}[1]{\left|#1\right|}
\newcommand{\norm}[1]{\left\lVert#1\right\rVert}

% 内积
\newcommand{\inner}[2]{\left\langle #1,#2\right\rangle}

% 常用算子
\DeclareMathOperator{\rank}{rank}
\DeclareMathOperator{\tr}{tr}
\DeclareMathOperator{\diag}{diag}
\DeclareMathOperator{\Span}{span}
\DeclareMathOperator{\supp}{supp}
\DeclareMathOperator{\Argmin}{arg\,min}
\DeclareMathOperator{\Argmax}{arg\,max}


% =========================================================
% 文档信息
% =========================================================

\title{数学知识讲义}
\subtitle{}

\author{}
\institute{}

\date{\today}

\version{1.0}

% 目录深度
\setcounter{tocdepth}{2}


% =========================================================
% 正文
% =========================================================

\begin{document}


% =========================================================
% 封面
% =========================================================

\maketitle


% =========================================================
% 前言部分
% =========================================================

\frontmatter

\tableofcontents


% =========================================================
% 正文部分
% =========================================================

\mainmatter


% =========================================================
\chapter{章节标题}
% =========================================================


% ---------------------------------------------------------
\begin{introduction}

    \item 
    \item 
    \item 
    \item 

\end{introduction}


% =========================================================
\section{问题背景}
% =========================================================



% =========================================================
\section{核心思想}
% =========================================================



% =========================================================
\section{基本定义}
% =========================================================

\begin{definition}[定义名称]



\end{definition}


% =========================================================
\section{直观理解}
% =========================================================

\begin{note}



\end{note}


% =========================================================
\section{几何意义}
% =========================================================



% =========================================================
\section{主要性质}
% =========================================================

\begin{proposition}[性质名称]



\end{proposition}


\begin{proof}



\end{proof}


% =========================================================
\section{重要定理}
% =========================================================

\begin{theorem}[定理名称]



\end{theorem}


\begin{proof}



\end{proof}


% =========================================================
\section{推论}
% =========================================================

\begin{corollary}[推论名称]



\end{corollary}


% =========================================================
\section{引理}
% =========================================================

\begin{lemma}[引理名称]



\end{lemma}


\begin{proof}



\end{proof}


% =========================================================
\section{典型例题}
% =========================================================

\begin{example}



\end{example}


\begin{solution}



\end{solution}


% =========================================================
\section{图像与几何解释}
% =========================================================

\begin{figure}[htbp]
    \centering

    % TikZ / PGFPlots 图像放置位置

    \caption{}
    \label{fig:example}
\end{figure}


% =========================================================
\section{知识联系}
% =========================================================



% =========================================================
\section{常见误区}
% =========================================================

\begin{remark}



\end{remark}


% =========================================================
\section{方法总结}
% =========================================================

\begin{conclusion}



\end{conclusion}


% =========================================================
\section{补充内容}
% =========================================================



% =========================================================
\chapter{练习}
% =========================================================

\begin{problemset}[练习题]

    \item

    \item

    \item

    \item

    \item

\end{problemset}


% =========================================================
% 附录
% =========================================================

\appendix


\chapter{补充公式}



\chapter{符号表}



% =========================================================
% 参考文献
% =========================================================

\backmatter

% 如果使用 BibTeX：
%
% \bibliographystyle{plain}
% \bibliography{references}


% =========================================================

\end{document}